%----------------------------------------------------------------------
\subsection{Calculation of the two-gluon incoming states}
\label{sec:wf-gg}
%----------------------------------------------------------------------

We now examine the quark contributions at order $g^{2}$ associated with an incoming state
containing two gluons, $\ket{\gl{a}{i}{k}\gl{b}{j}{p}}$. The allowed final configurations are a
quark--antiquark pair (virtual corrections to the pair sector and $gg\to q\bar q$ scattering),
a pair with two gluons, and two pairs. The disconnected configurations, in which one of the
incoming gluons is a spectator, reproduce the coefficients $\mf A_{3}$, $\mf A_{1}$ and $\mf B_{3}$
already computed and are not repeated.

\subsubsection{Quark--antiquark outgoing}

The overall energy denominator is
\begin{equation}
D_{f}\equiv\eps_{k}+\eps_{p}-\eps_{l}-\eps_{m}.
\end{equation}
Six topologies contribute.

\paragraph{(a) $s$-channel: the two gluons merge and the resulting gluon splits.}
\begin{center}
\begin{tikzpicture}
\draw[gl] (0,1.9) -- (1.4,1.25); \node[left] at (0,1.9) {\scriptsize$\gl{a}{i}{k}$};
\draw[gl] (0,0.6) -- (1.4,1.25); \node[left] at (0,0.6) {\scriptsize$\gl{b}{j}{p}$};
\draw[gl] (1.4,1.25) -- (2.5,1.25); \node[above] at (1.95,1.3) {\scriptsize$\gl{c'}{k'}{q'}$};
\draw[qk] (2.5,1.25) -- (3.9,1.85); \node[right] at (3.9,1.85) {\scriptsize$\qk{\beta}{\lambda_{1}}{l}$};
\draw[qk] (3.9,0.55) -- (2.5,1.25); \node[right] at (3.9,0.55) {\scriptsize$\aqk{\alpha}{\lambda_{2}}{m}$};
\draw[cut] (1.95,0.1) -- (1.95,2.25); \draw[cut] (3.2,0.1) -- (3.2,2.25);
\end{tikzpicture}
\end{center}
\begin{equation}
\ket{(\Psi^{ab}_{ij}(k,p))^{1}_{q\bar q}}=\int\! d^{3}q'\,d^{3}l\,d^{3}m\;
\ket{q\bar q}\;
\frac{\bra{q\bar q}H_{gq\bar q}\ket{\gl{c'}{k'}{q'}}}{D_{f}}\;
\frac{\bra{\gl{c'}{k'}{q'}}H_{ggg}\ket{\gl{a}{i}{k}\gl{b}{j}{p}}}{\eps_{k}+\eps_{p}-\eps_{q'}} .
\end{equation}
The $2\to1$ gluon vertex is the Hermitian conjugate of \eqref{eq:MEggg}; with $q'=k+p$,
\begin{multline}
\ket{(\Psi^{ab}_{ij}(k,p))^{1}_{q\bar q}}=-\frac{ig^{2}}{64\pi^{3}}\int\!d^{3}l\,d^{3}m\;
\\
\times\ket{q\bar q}\,f^{abc'}t^{c'}_{\beta\alpha}\,
\frac{\dthree(k+p-l-m)}{q'^{+}\sqrt{k^{+}p^{+}}}\,
\frac{\Vg^{k'ij}(q';k,p)\;\Ph^{k'}_{\lambda_{1}\lambda_{2}}(q';l,m)}
{\big(\eps_{k}+\eps_{p}-\eps_{q'}\big)\,D_{f}} .
\label{eq:B1a}
\end{multline}
This contribution is symmetric under $(a,i,k)\leftrightarrow(b,j,p)$.

\paragraph{(b) One gluon is absorbed by the valence charge, the other splits.}
\begin{center}
\begin{tikzpicture}
\draw[val] (0,0) -- (3.9,0); \node[left] at (-0.05,0) {\scriptsize$\rho$};
\draw[gl] (0,0.9) -- (2.9,0); \node[left] at (0,0.9) {\scriptsize$\gl{a}{i}{k}$};
\draw[gl] (0,2.1) -- (1.3,2.1); \node[left] at (0,2.1) {\scriptsize$\gl{b}{j}{p}$};
\draw[qk] (1.3,2.1) -- (3.9,2.5); \node[right] at (3.9,2.5) {\scriptsize$\qk{\beta}{\lambda_{1}}{l}$};
\draw[qk] (3.9,1.5) -- (1.3,2.1); \node[right] at (3.9,1.5) {\scriptsize$\aqk{\alpha}{\lambda_{2}}{m}$};
\node at (1.95,-0.55) {\scriptsize (b1)};
\end{tikzpicture}
\hspace{0.5cm}
\begin{tikzpicture}
\draw[val] (0,0) -- (3.9,0); \node[left] at (-0.05,0) {\scriptsize$\rho$};
\draw[gl] (0,0.9) -- (1.1,0); \node[left] at (0,0.9) {\scriptsize$\gl{a}{i}{k}$};
\draw[gl] (0,2.1) -- (2.6,2.1); \node[left] at (0,2.1) {\scriptsize$\gl{b}{j}{p}$};
\draw[qk] (2.6,2.1) -- (3.9,2.5); \node[right] at (3.9,2.5) {\scriptsize$\qk{\beta}{\lambda_{1}}{l}$};
\draw[qk] (3.9,1.5) -- (2.6,2.1); \node[right] at (3.9,1.5) {\scriptsize$\aqk{\alpha}{\lambda_{2}}{m}$};
\node at (1.95,-0.55) {\scriptsize (b2)};
\end{tikzpicture}
\end{center}
In (b1) the intermediate state is $\ket{q\bar q\,\gl{a}{i}{k}}$, with
$E_{i}-E_{m}=\eps_{p}-\eps_{l}-\eps_{m}$; in (b2) it is $\ket{\gl{b}{j}{p}}$, with $E_{i}-E_{m}=\eps_{k}$:
\begin{align}
\ket{(\Psi^{ab}_{ij}(k,p))^{2}_{q\bar q\,\rho}}&=\int\!d^{3}l\,d^{3}m\;\ket{q\bar q}\;
\frac{\bra{q\bar q}H_{g}\ket{q\bar q\,\gl{a}{i}{k}}}{D_{f}}\;
\frac{\bra{q\bar q\,\gl{a}{i}{k}}H_{gq\bar q}\ket{\gl{a}{i}{k}\gl{b}{j}{p}}}{\eps_{p}-\eps_{l}-\eps_{m}},
\\
\ket{(\Psi^{ab}_{ij}(k,p))^{3}_{q\bar q\,\rho}}&=\int\!d^{3}l\,d^{3}m\;\ket{q\bar q}\;
\frac{\bra{q\bar q}H_{gq\bar q}\ket{\gl{b}{j}{p}}}{D_{f}}\;
\frac{\bra{\gl{b}{j}{p}}H_{g}\ket{\gl{a}{i}{k}\gl{b}{j}{p}}}{\eps_{k}} .
\end{align}
Since $D_{f}=(\eps_{p}-\eps_{l}-\eps_{m})+\eps_{k}$, the two orderings combine by \eqref{eq:factor},
\begin{multline}
\ket{(\Psi^{ab}_{ij}(k,p))^{2+3}_{q\bar q\,\rho}}=\frac{g^{2}}{32\pi^{3}}\int\!d^{3}l\,d^{3}m\;
\ket{q\bar q}\;
\frac{\rho^{a}(k)\,t^{b}_{\beta\alpha}\,k^{i}\,\Ph^{j}_{\lambda_{1}\lambda_{2}}(p;l,m)\,\dthree(p-l-m)}
{(k^{+})^{3/2}\sqrt{p^{+}}\;\eps_{k}\,\big(\eps_{p}-\eps_{l}-\eps_{m}\big)}\\
+\big((a,i,k)\leftrightarrow(b,j,p)\big) .
\label{eq:B1b}
\end{multline}

\paragraph{(c) $t$-channel quark exchange: one gluon splits and the quark absorbs the other.}
\begin{center}
\begin{tikzpicture}
\draw[gl] (0,1.5) -- (1.2,1.5); \node[left] at (0,1.5) {\scriptsize$\gl{a}{i}{k}$};
\draw[qk] (1.2,1.5) -- (2.5,2.05); \node[above left] at (1.95,1.75) {\scriptsize$l_{1}$};
\draw[qk] (2.5,2.05) -- (3.9,2.3); \node[right] at (3.9,2.3) {\scriptsize$\qk{\beta}{\lambda_{1}}{l}$};
\draw[qk] (3.9,0.9) -- (1.2,1.5); \node[right] at (3.9,0.9) {\scriptsize$\aqk{\alpha}{\lambda_{2}}{m}$};
\draw[gl] (0,2.6) -- (2.5,2.05); \node[left] at (0,2.6) {\scriptsize$\gl{b}{j}{p}$};
\draw[cut] (1.85,0.5) -- (1.85,2.9); \draw[cut] (3.2,0.5) -- (3.2,2.9);
\node at (1.95,0.1) {\scriptsize (c)};
\end{tikzpicture}
\hspace{0.8cm}
\begin{tikzpicture}
\draw[gl] (0,1.5) -- (1.2,1.5); \node[left] at (0,1.5) {\scriptsize$\gl{a}{i}{k}$};
\draw[qk] (1.2,1.5) -- (3.9,2.3); \node[right] at (3.9,2.3) {\scriptsize$\qk{\beta}{\lambda_{1}}{l}$};
\draw[qk] (2.5,0.95) -- (1.2,1.5); \node[below left] at (1.95,1.1) {\scriptsize$m_{1}$};
\draw[qk] (3.9,0.7) -- (2.5,0.95); \node[right] at (3.9,0.7) {\scriptsize$\aqk{\alpha}{\lambda_{2}}{m}$};
\draw[gl] (0,0.3) -- (2.5,0.95); \node[left] at (0,0.3) {\scriptsize$\gl{b}{j}{p}$};
\draw[cut] (1.85,0.0) -- (1.85,2.6); \draw[cut] (3.2,0.0) -- (3.2,2.6);
\node at (1.95,-0.4) {\scriptsize (d)};
\end{tikzpicture}
\end{center}
The gluon $k$ splits into $\qk{\gamma}{\lambda}{l_{1}}\aqk{\alpha}{\lambda_{2}}{m}$ and the quark absorbs
the gluon $p$, so that $l_{1}=l-p$ and the process requires $l^{+}>p^{+}$. The absorption vertex
is the Hermitian conjugate of \eqref{eq:MEqg}, $\big[\Ph^{j}_{\lambda'\lambda}(p;l_{1},l)\big]^{*}
=\Ph^{j}_{\lambda\lambda'}(p;l,l_{1})$:
\begin{align}
\ket{(\Psi^{ab}_{ij}(k,p))^{4}_{q\bar q}}&=\int\!d^{3}l_{1}\,d^{3}l\,d^{3}m\;\ket{q\bar q}\;
\frac{\bra{\qk{\beta}{\lambda_{1}}{l}}H_{gq\bar q}\ket{\qk{\gamma}{\lambda}{l_{1}}\gl{b}{j}{p}}}{D_{f}}
\nonumber\\&\hspace{3.2cm}\times
\frac{\bra{\qk{\gamma}{\lambda}{l_{1}}\aqk{\alpha}{\lambda_{2}}{m}}H_{gq\bar q}\ket{\gl{a}{i}{k}}}
{\eps_{k}-\eps_{l_{1}}-\eps_{m}}
\nonumber\\
&=\frac{g^{2}}{64\pi^{3}}\int\!d^{3}l\,d^{3}m\;\ket{q\bar q}\;(t^{b}t^{a})_{\beta\alpha}\,
\frac{\dthree(k+p-l-m)\,\theta(l^{+}-p^{+})}{\sqrt{k^{+}p^{+}}}
\nonumber\\&\hspace{3.2cm}\times
\frac{\chi^{\dg}_{\lambda_{1}}\Gam^{j}(p;l,l_{1})\,\Gam^{i}(k;l_{1},m)\chi_{\lambda_{2}}}
{\big(\eps_{k}-\eps_{l_{1}}-\eps_{m}\big)\,D_{f}}
\nonumber\\
&\quad+\big((a,i,k)\leftrightarrow(b,j,p)\big).
\label{eq:B1c}
\end{align}

\paragraph{(d) $u$-channel quark exchange: one gluon splits and the antiquark absorbs the other.}
Now $m_{1}=m-p$, the process requires $m^{+}>p^{+}$, and the antiquark vertex \eqref{eq:MEaqg}
supplies a minus sign:
\begin{multline}
\ket{(\Psi^{ab}_{ij}(k,p))^{5}_{q\bar q}}=-\frac{g^{2}}{64\pi^{3}}\int\!d^{3}l\,d^{3}m\;\ket{q\bar q}\;\\
\times(t^{a}t^{b})_{\beta\alpha}\,
\frac{\dthree(k+p-l-m)\,\theta(m^{+}-p^{+})}{\sqrt{k^{+}p^{+}}}\,
\frac{\chi^{\dg}_{\lambda_{1}}\Gam^{i}(k;l,m_{1})\,\Gam^{j}(p;m_{1},m)\chi_{\lambda_{2}}}
{\big(\eps_{k}-\eps_{l}-\eps_{m_{1}}\big)\,D_{f}}\\
+\big((a,i,k)\leftrightarrow(b,j,p)\big).
\label{eq:B1d}
\end{multline}
The diagrams (c) and (d) are the light-cone time ordered $t$- and $u$-channel quark exchange
graphs of $gg\to q\bar q$. In the language of the evolution operator they are the products
$\mf F_{2}\,\mf A_{3}$ and $\mf G_{2}\,\mf A_{3}$ in \eqref{eq:reduction}, and they are the origin of
the $\mf F$ and $\mf G$ terms of the unitarity relation \eqref{eq:UB2}.

\paragraph{(e) Instantaneous quark exchange.}
\begin{center}
\begin{tikzpicture}
\draw[gl] (0,2.0) -- (1.3,1.35); \node[left] at (0,2.0) {\scriptsize$\gl{a}{i}{k}$};
\draw[gl] (0,0.1) -- (2.3,0.8); \node[left] at (0,0.1) {\scriptsize$\gl{b}{j}{p}$};
\draw[qk] (1.3,1.35) -- (3.9,1.9); \node[right] at (3.9,1.9) {\scriptsize$\qk{\beta}{\lambda_{1}}{l}$};
\draw[ins] (1.3,1.35) -- (2.3,0.8); \ibarh{1.8}{1.075}
\draw[qk] (3.9,0.25) -- (2.3,0.8); \node[right] at (3.9,0.25) {\scriptsize$\aqk{\alpha}{\lambda_{2}}{m}$};
\draw[cut] (3.2,-0.15) -- (3.2,2.2);
\end{tikzpicture}
\end{center}
With \eqref{eq:MEggqq2},
\begin{multline}
\ket{(\Psi^{ab}_{ij}(k,p))^{6}_{q\bar q}}=\frac{g^{2}}{16\pi^{3}}\int\!d^{3}l\,d^{3}m\;\ket{q\bar q}\,
\\
\times\frac{\dthree(k+p-l-m)}{\sqrt{k^{+}p^{+}}\,D_{f}}\,
\chi^{\dg}_{\lambda_{1}}\!\left[\frac{(t^{a}t^{b})_{\beta\alpha}\sigma^{i}\sigma^{j}}{l^{+}-k^{+}}
+\frac{(t^{b}t^{a})_{\beta\alpha}\sigma^{j}\sigma^{i}}{l^{+}-p^{+}}\right]\chi_{\lambda_{2}} .
\label{eq:B1e}
\end{multline}
The second term is the instantaneous part of the $t$-channel exchange (c) (the fermion line
carries $l_{1}^{+}=l^{+}-p^{+}$) and the first that of the $u$-channel exchange (d) (it carries
$l^{+}-k^{+}=-m_{1}^{+}$); unlike (c) and (d), the instantaneous term has support for either sign
of these momenta.

\paragraph{(f) Instantaneous gluon exchange between the gluon and quark currents.}
\begin{center}
\begin{tikzpicture}
\draw[gl] (0,2.0) -- (1.8,1.5); \node[left] at (0,2.0) {\scriptsize$\gl{a}{i}{k}$};
\draw[gl] (0,1.0) -- (1.8,1.5); \node[left] at (0,1.0) {\scriptsize$\gl{b}{j}{p}$};
\draw[ins] (1.8,1.5) -- (1.8,0.6); \ibar{1.8}{1.05}
\draw[qk] (1.8,0.6) -- (3.9,1.05); \node[right] at (3.9,1.05) {\scriptsize$\qk{\beta}{\lambda_{1}}{l}$};
\draw[qk] (3.9,0.0) -- (1.8,0.6); \node[right] at (3.9,0.0) {\scriptsize$\aqk{\alpha}{\lambda_{2}}{m}$};
\draw[cut] (3.2,-0.25) -- (3.2,2.2);
\end{tikzpicture}
\end{center}
The two incoming gluons are annihilated by $J_{g}$ and the pair is created by $J_{q}$; the
exchanged gluon carries $P^{+}=k^{+}+p^{+}$ and, with \eqref{eq:MEJJ2},
\begin{equation}
\ket{(\Psi^{ab}_{ij}(k,p))^{7}_{q\bar q}}=-\frac{ig^{2}}{2(2\pi)^{3}}\int\!d^{3}l\,d^{3}m\;\ket{q\bar q}\,
\frac{f^{cab}\,t^{c}_{\beta\alpha}\,\delta_{\lambda_{1}\lambda_{2}}\,\delta_{ij}\,(p^{+}-k^{+})\,
\dthree(k+p-l-m)}{\sqrt{k^{+}p^{+}}\,(k^{+}+p^{+})^{2}\,D_{f}} .
\label{eq:B1f}
\end{equation}

\subsubsection{Two gluons, quark and antiquark outgoing}

One incoming gluon splits into the pair and the other into the two outgoing gluons
$\gl{d}{o}{r}\gl{e}{n}{s}$; both time orderings contribute.
\begin{center}
\begin{tikzpicture}
\draw[gl] (0,2.2) -- (1.4,2.2); \node[left] at (0,2.2) {\scriptsize$\gl{a}{i}{k}$};
\draw[qk] (1.4,2.2) -- (3.9,2.6); \node[right] at (3.9,2.6) {\scriptsize$\qk{\beta}{\lambda_{1}}{l}$};
\draw[qk] (3.9,1.7) -- (1.4,2.2); \node[right] at (3.9,1.7) {\scriptsize$\aqk{\alpha}{\lambda_{2}}{m}$};
\draw[gl] (0,0.5) -- (2.2,0.5); \node[left] at (0,0.5) {\scriptsize$\gl{b}{j}{p}$};
\draw[gl] (2.2,0.5) -- (3.9,0.95); \node[right] at (3.9,0.95) {\scriptsize$\gl{d}{o}{r}$};
\draw[gl] (2.2,0.5) -- (3.9,0.0); \node[right] at (3.9,0.0) {\scriptsize$\gl{e}{n}{s}$};
\draw[cut] (1.8,-0.35) -- (1.8,2.95); \draw[cut] (3.2,-0.35) -- (3.2,2.95);
\end{tikzpicture}
\end{center}
With $D_{1}=\eps_{k}-\eps_{l}-\eps_{m}$, $D_{2}=\eps_{p}-\eps_{r}-\eps_{s}$ and $D_{f}=D_{1}+D_{2}$,
\eqref{eq:factor} gives
\begin{align}
\ket{(\Psi^{ab}_{ij}(k,p))_{ggq\bar q}}=\;&\frac{ig^{2}}{64\pi^{3}}\int\!d^{3}l\,d^{3}m\,d^{3}r\,d^{3}s\;
\ket{q\bar q\,\gl{d}{o}{r}\gl{e}{n}{s}}
\nonumber\\&\times t^{a}_{\beta\alpha}f^{bde}\,
\frac{\dthree(k-l-m)\dthree(p-r-s)}{\sqrt{k^{+}p^{+}r^{+}s^{+}}}\,
\frac{\Ph^{i}_{\lambda_{1}\lambda_{2}}(k;l,m)\;\Vg^{jon}(p;r,s)}{D_{1}\,D_{2}}
\nonumber\\&+\big((a,i,k)\leftrightarrow(b,j,p)\big),
\label{eq:Pggqq}
\end{align}
where the factor $\tfrac12$ of the two identical outgoing gluons has cancelled against their two
assignments, as in Paper~I.

\subsubsection{Two quark--antiquark pairs outgoing}

Each gluon splits independently and both orderings again combine by \eqref{eq:factor}:
\begin{align}
\ket{(\Psi^{bc}_{jk}(p,q))_{q\bar qq\bar q}}=\frac{g^{2}}{64\pi^{3}}
\int\!d^{3}l\,d^{3}m\,d^{3}l'\,d^{3}m'\;
&\ket{q\bar q\,q\bar q}\;t^{b}_{\beta\alpha}t^{c}_{\beta'\alpha'}\,
\frac{\dthree(p-l-m)\dthree(q-l'-m')}{\sqrt{p^{+}q^{+}}}
\nonumber\\&\times
\frac{\Ph^{j}_{\lambda_{1}\lambda_{2}}(p;l,m)\;\Ph^{k}_{\lambda_{1}'\lambda_{2}'}(q;l',m')}
{\big(\eps_{p}-\eps_{l}-\eps_{m}\big)\big(\eps_{q}-\eps_{l'}-\eps_{m'}\big)} .
\label{eq:Pqqqq}
\end{align}

\subsubsection{Summary of results}
\label{sec:sum-gg}

The connected quark components of the two-gluon state at order $g^{2}$ are
\begin{multline}
\ket{\Psi^{ab}_{ij}(k,p)}\supset\sum_{n=1,4,5,6,7}\ket{(\Psi^{ab}_{ij}(k,p))^{n}_{q\bar q}}
+\ket{(\Psi^{ab}_{ij}(k,p))^{2+3}_{q\bar q\,\rho}}\\
+\ket{(\Psi^{ab}_{ij}(k,p))_{ggq\bar q}}+\ket{(\Psi^{ab}_{ij}(k,p))_{q\bar qq\bar q}} .
\label{eq:sum-gg}
\end{multline}
The $\rho$ dependent component and the components with four outgoing partons factorize into
products of first order amplitudes. The $\rho$ independent part of the $q\bar q$ component is the
$gg\to q\bar q$ scattering amplitude; its overall denominator $D_{f}$ vanishes on the energy
shell and is understood with the prescription of section~\ref{sec:lcpt}. The new connected
components do not affect the normalization at this order.

%----------------------------------------------------------------------
\subsection{Calculation of the three-gluon incoming states}
\label{sec:wf-ggg}
%----------------------------------------------------------------------

As in Paper~I, we restrict the calculation to the components needed for $\OM$ at
$\order{g^{2}}$. For an incoming state $\ket{\gl{a}{i}{k}\gl{b}{j}{p}\gl{c}{k}{q}}$ the only new
connected quark component is $\ket{q\bar q\,\gl{e}{n}{s}}$: one gluon splits into the pair while
the other two merge.
\begin{center}
\begin{tikzpicture}
\draw[gl] (0,2.3) -- (1.5,2.3); \node[left] at (0,2.3) {\scriptsize$\gl{a}{i}{k}$};
\draw[qk] (1.5,2.3) -- (3.9,2.7); \node[right] at (3.9,2.7) {\scriptsize$\qk{\beta}{\lambda_{1}}{l}$};
\draw[qk] (3.9,1.8) -- (1.5,2.3); \node[right] at (3.9,1.8) {\scriptsize$\aqk{\alpha}{\lambda_{2}}{m}$};
\draw[gl] (0,0.9) -- (2.3,0.45); \node[left] at (0,0.9) {\scriptsize$\gl{b}{j}{p}$};
\draw[gl] (0,0.0) -- (2.3,0.45); \node[left] at (0,0.0) {\scriptsize$\gl{c}{k}{q}$};
\draw[gl] (2.3,0.45) -- (3.9,0.45); \node[right] at (3.9,0.45) {\scriptsize$\gl{e}{n}{s}$};
\draw[cut] (1.9,-0.35) -- (1.9,3.05); \draw[cut] (3.2,-0.35) -- (3.2,3.05);
\end{tikzpicture}
\end{center}
With $D_{A}=\eps_{k}-\eps_{l}-\eps_{m}$ and $D_{B}=\eps_{p}+\eps_{q}-\eps_{s}$, $D_{A}+D_{B}=D_{f}$,
and \eqref{eq:factor} gives
\begin{align}
\ket{(\Psi^{abc}_{ijk}(k,p,q))_{gq\bar q}}=-\frac{ig^{2}}{64\pi^{3}}
\int\!d^{3}l\,d^{3}m\,d^{3}s\;
&\ket{q\bar q\,\gl{e}{n}{s}}\;t^{a}_{\beta\alpha}f^{ebc}\,
\frac{\dthree(k-l-m)\dthree(s-p-q)}{\sqrt{k^{+}s^{+}p^{+}q^{+}}}
\nonumber\\&\times
\frac{\Ph^{i}_{\lambda_{1}\lambda_{2}}(k;l,m)\;\Vg^{njk}(s;p,q)}{D_{A}\,D_{B}} ,
\label{eq:Pgqq}
\end{align}
summed over the three inequivalent choices of the gluon that splits. Since
$f^{ebc}\Vg^{njk}(s;p,q)$ is symmetric under $(b,j,p)\leftrightarrow(c,k,q)$, a sum over all six
permutations would double count.

\subsubsection{Summary of results}

The only new connected quark component of the three-gluon state at order $g^{2}$ is
\eqref{eq:Pgqq}; it factorizes into the product of the pair splitting amplitude and the
three-gluon merging amplitude of Paper~I.
